% GENERATED FILE. Edit canonical lessons, then run npm run generate:books.
% canonical-source-hash: fnv1a64:23d47d043175c8f3
% canonical-lessons: KA-C236-biccu, KA-C236-ebbisu, KA-C236-eccaragollu, KA-C236-caukasi-madu, KA-C236-eru

\chapter{To untie, To wake, To wake up, To bargain, To climb}
\label{ch:ka-to-untie-to-wake-to-wake-up-to-bargain-to-climb}
\hlchaptermodality{236}

\begin{chapteropening}
\textbf{By the end of this chapter:} \emph{I can say to untie, to wake, to wake up, to bargain and to climb.}

Complete the last lesson of chapter 236: I can say to untie, to wake, to wake up, to bargain and to climb.
\end{chapteropening}

\section[\texorpdfstring{\kn{ಬಿಚ್ಚು} (biccu)}{ಬಿಚ್ಚು (biccu)}]{\kn{ಬಿಚ್ಚು} (biccu) — to untie, to take off (clothes)}

\label{lesson:KA-C236-biccu}



\begin{quote}
Before the new one: say the Kannada for to lift, then the Kannada for to throw.
\end{quote}

\subsection*{You'll want to know: \kn{ಬಿಚ್ಚು}}
\textbf{\kn{ಬಿಚ್ಚು}} — \emph{biccu} — \textquotedblleft{}to untie, to take off (clothes)\textquotedblright{}.

\begin{grammarlens}[title={What you've built}]
One more word for this chapter.
\end{grammarlens}

\subsection*{Your turn}
\begin{itemize}
\raggedright
  \item \emph{Say it:} \emph{biccu}
  \item \emph{Say it:} \emph{biccu}, once more
  \item \emph{Say it:} say \emph{biccu}
  \item \emph{Recall:} say the Kannada for to lift, then the Kannada for to throw, then say \emph{biccu} again
\end{itemize}

\begin{tcolorbox}[breakable,colback=teal!4,colframe=teal!35!black,title={Before you move on}]
What does \kn{ಬಿಚ್ಚು} mean? (\textquotedblleft{}To untie, to take off (clothes)\textquotedblright{}.) Say it once more.
\end{tcolorbox}
\section[\texorpdfstring{\kn{ಎಬ್ಬಿಸು} (ebbisu)}{ಎಬ್ಬಿಸು (ebbisu)}]{\kn{ಎಬ್ಬಿಸು} (ebbisu) — to wake (someone) up}

\label{lesson:KA-C236-ebbisu}



\begin{quote}
Before the new one: say the Kannada for to throw, then the Kannada for to untie.
\end{quote}

\subsection*{You'll want to know: \kn{ಎಬ್ಬಿಸು}}
\textbf{\kn{ಎಬ್ಬಿಸು}} — \emph{ebbisu} — \textquotedblleft{}to wake (someone) up\textquotedblright{}.

\begin{grammarlens}[title={What you've built}]
One more word for this chapter.
\end{grammarlens}

\subsection*{Your turn}
\begin{itemize}
\raggedright
  \item \emph{Say it:} \emph{ebbisu}
  \item \emph{Say it:} \emph{ebbisu}, once more
  \item \emph{Say it:} say \emph{ebbisu}
  \item \emph{Recall:} say the Kannada for to throw, then the Kannada for to untie, then say \emph{ebbisu} again
\end{itemize}

\begin{tcolorbox}[breakable,colback=teal!4,colframe=teal!35!black,title={Before you move on}]
What does \kn{ಎಬ್ಬಿಸು} mean? (\textquotedblleft{}To wake (someone) up\textquotedblright{}.) Say it once more.
\end{tcolorbox}
\section[\texorpdfstring{\kn{ಎಚ್ಚರಗೊಳ್ಳು} (eccaragoḷḷu)}{ಎಚ್ಚರಗೊಳ್ಳು (eccaragoḷḷu)}]{\kn{ಎಚ್ಚರಗೊಳ್ಳು} (eccaragoḷḷu) — to wake up}

\label{lesson:KA-C236-eccaragollu}



\begin{quote}
Before the new one: say the Kannada for to untie, then the Kannada for to wake.
\end{quote}

\subsection*{You'll want to know: \kn{ಎಚ್ಚರಗೊಳ್ಳು}}
\textbf{\kn{ಎಚ್ಚರಗೊಳ್ಳು}} — \emph{eccaragoḷḷu} — \textquotedblleft{}to wake up\textquotedblright{}.

\begin{grammarlens}[title={What you've built}]
One more word for this chapter.
\end{grammarlens}

\subsection*{Your turn}
\begin{itemize}
\raggedright
  \item \emph{Say it:} \emph{eccaragoḷḷu}
  \item \emph{Say it:} \emph{eccaragoḷḷu}, once more
  \item \emph{Say it:} say \emph{eccaragoḷḷu}
  \item \emph{Recall:} say the Kannada for to untie, then the Kannada for to wake, then say \emph{eccaragoḷḷu} again
\end{itemize}

\begin{tcolorbox}[breakable,colback=teal!4,colframe=teal!35!black,title={Before you move on}]
What does \kn{ಎಚ್ಚರಗೊಳ್ಳು} mean? (\textquotedblleft{}To wake up\textquotedblright{}.) Say it once more.
\end{tcolorbox}
\section[\texorpdfstring{\kn{ಚೌಕಾಸಿ} \kn{ಮಾಡು} (caukāsi māḍu)}{ಚೌಕಾಸಿ ಮಾಡು (caukāsi māḍu)}]{\kn{ಚೌಕಾಸಿ} \kn{ಮಾಡು} (caukāsi māḍu) — to bargain}

\label{lesson:KA-C236-caukasi-madu}



\begin{quote}
Before the new one: say the Kannada for to wake, then the Kannada for to wake up.
\end{quote}

\subsection*{You'll want to know: \kn{ಚೌಕಾಸಿ} \kn{ಮಾಡು}}
\textbf{\kn{ಚೌಕಾಸಿ} \kn{ಮಾಡು}} — \emph{caukāsi māḍu} — \textquotedblleft{}to bargain\textquotedblright{}.

\begin{grammarlens}[title={What you've built}]
One more word for this chapter.
\end{grammarlens}

\subsection*{Your turn}
\begin{itemize}
\raggedright
  \item \emph{Say it:} \emph{caukāsi māḍu}
  \item \emph{Say it:} \emph{caukāsi māḍu}, once more
  \item \emph{Say it:} say \emph{caukāsi māḍu}
  \item \emph{Recall:} say the Kannada for to wake, then the Kannada for to wake up, then say \emph{caukāsi māḍu} again
\end{itemize}

\begin{tcolorbox}[breakable,colback=teal!4,colframe=teal!35!black,title={Before you move on}]
What does \kn{ಚೌಕಾಸಿ} \kn{ಮಾಡು} mean? (\textquotedblleft{}To bargain\textquotedblright{}.) Say it once more.
\end{tcolorbox}
\section[\texorpdfstring{\kn{ಏರು} (ēru)}{ಏರು (ēru)}]{\kn{ಏರು} (ēru) — to climb, to go up}

\label{lesson:KA-C236-eru}



\begin{quote}
Before the new one: say the Kannada for to wake up, then the Kannada for to bargain.
\end{quote}

\subsection*{You'll want to know: \kn{ಏರು}}
\textbf{\kn{ಏರು}} — \emph{ēru} — \textquotedblleft{}to climb, to go up\textquotedblright{}.

\begin{grammarlens}[title={What you've built}]
One more word for this chapter.
\end{grammarlens}

\subsection*{Your turn}
\begin{itemize}
\raggedright
  \item \emph{Say it:} \emph{ēru}
  \item \emph{Say it:} \emph{ēru}, once more
  \item \emph{Say it:} say \emph{ēru}
  \item \emph{Recall:} say the Kannada for to wake up, then the Kannada for to bargain, then say \emph{ēru} again
\end{itemize}

\begin{tcolorbox}[breakable,colback=teal!4,colframe=teal!35!black,title={Before you move on}]
What does \kn{ಏರು} mean? (\textquotedblleft{}To climb, to go up\textquotedblright{}.) Say it once more.
\end{tcolorbox}
